% Source: https://artofproblemsolving.com/community/c3h1919179_geometry_question
% Stable identity: aops-topic-1919179-q01
From the ratios,
\[
AX=\dfrac{3}{4}AB,\ XB=\dfrac{1}{4}AB,\qquad BY=\dfrac{4}{5}BC,\ YC=\dfrac{1}{5}BC,\qquad CZ=\dfrac{5}{6}CA,\ ZA=\dfrac{1}{6}CA .
\]
Two triangles with a common angle have areas in the ratio of the products of the sides enclosing that angle (area $=\dfrac{1}{2}\cdot$ side $\cdot$ side $\cdot\sin$ angle). Hence
\[
\dfrac{[AXZ]}{[ABC]}=\dfrac{AX\cdot AZ}{AB\cdot AC}=\dfrac{3}{4}\cdot\dfrac{1}{6}=\dfrac{1}{8},\quad
\dfrac{[BYX]}{[ABC]}=\dfrac{BX\cdot BY}{BA\cdot BC}=\dfrac{1}{4}\cdot\dfrac{4}{5}=\dfrac{1}{5},\quad
\dfrac{[CZY]}{[ABC]}=\dfrac{CY\cdot CZ}{CB\cdot CA}=\dfrac{1}{5}\cdot\dfrac{5}{6}=\dfrac{1}{6} .
\]
So the three corner triangles have areas $30$, $48$ and $40$:

\begin{asy}
usepackage("amsmath");
size(6.5cm);
defaultpen(fontsize(10pt));
pair A=(0,0),B=(10,0),C=(3.5,7);
pair X=A+0.75*(B-A), Y=B+0.8*(C-B), Z=C+(5/6)*(A-C);
fill(X--Y--Z--cycle,mediumgray+white);
draw(A--B--C--cycle,linewidth(1)); draw(X--Y--Z--cycle,linewidth(0.8));
label("$A$",A,SW);label("$B$",B,SE);label("$C$",C,N);
label("$X$",X,S);label("$Y$",Y,NE);label("$Z$",Z,W);
label("$30$",(2.6,0.38));label("$48$",(7.75,1.2));label("$40$",(C+Y+Z)/3);label("$122$",(X+Y+Z)/3);
label("\mbox{}",(min(currentpicture,user=true)+max(currentpicture,user=true))/2,black); // keeps the text after the figure black when LaTeX typesets the labels inline
\end{asy}

and
\[
[XYZ]=240-30-48-40=\boxed{122}.
\]
